% Page 44 — suite de « Alban : XVIème siècle (début) »

seraient les fondateurs et celle de Sainte Anne fondée en 1477 par Paulin Limouzin (ou Limouzi) le 28 janvier de cette année-là. La susdite famille des Lemouzy est encore à Alban en 1598, l'un d'eux Meric Lemouzy possède une terre à La Bessière (là où sont actuellement les Bascoul) confrontant le chemin d'Alban à Réalmont et le chemin d'Alban à la Berrière aux terres de Jean Valette.

\medskip

\lettrine{À}{ la} même époque, disparaît l'abbaye Saint Benoît de Castres que remplace en 533 le chapitre cathédral de la même ville. Et c'est à ce même chapitre cathédral que sera rattachée, au moins juridiquement, la paroisse de Notre Dame d'Alban.

\medskip

\lettrine{D}{e}puis le dernier pillage de 1484 qui a suivi la rébellion des Albanais contre le délégué royal Boffille de Juge, jusqu'à 1562, c'est presque quatre vingt ans de paix dont jouira notre pays. Est-ce en ce temps-là qu'Alban a pu compter jusqu'à 400 feux (soit au bas mot 2000 habitants) ? Le cahier des doléances de 1789 y fera allusion. Mais nos délégués municipaux n'ont-ils pas gonflé un peu trop le chiffre, un peu à la manière des Maires d'aujourd'hui, afin d'obtenir des subventions désirables ? La chose n'est pas invraisemblable. D'autant que le registre cadastral de 1598 ne fera état que d'un chiffre inférieur à soixante feux (ou familles).
